\chapter{反应为什么会放热或吸热}

\begin{kaichang}
冬天的教室里，总有同学从口袋里摸出一片暖宝宝：撕开塑料袋，搓一搓，几分钟后它就开始发热，能暖暖和和地焐上小半天。没有火，没有电，没有电池——热是从哪儿来的？

医院和运动场上还有另一种相反的小袋子：一次性冰袋。用力一捏，里面“啪”地一声，某个内袋破了，几秒钟后整个袋子变得冰凉，能敷在扭伤的脚踝上。没有冰箱，没有压缩机——冷又是从哪儿来的？

一个凭空发热，一个凭空变凉。先猜一猜：暖宝宝里的热，是原本就“藏”在袋子里的吗？冰袋的凉，是它“制造”出来的吗？把答案写在纸上。这一章结束时，我们要回来给这两只小袋子算一笔完整的能量账。
\end{kaichang}

\section{先分清：谁是系统，谁是环境}

把暖宝宝攥在手里，你感觉到的是“它在发热”。但严格地说，你感觉到的其实是：\textbf{能量正在从袋子里流进你的手}。这个区分听起来抠字眼，却是整章的地基。

化学家把要研究的那一小部分世界叫作\textbf{系统}——比如暖宝宝里的铁粉和空气，或者烧杯里正在反应的溶液；系统之外的一切叫作\textbf{环境}——你的手、烧杯、房间里的空气。化学反应发生时，如果系统的总能量下降，多余的能量就以热的形式流向环境，环境的温度升高，我们就说这个反应\textbf{放热}；反过来，如果系统的总能量升高，它就得从环境“抽血”，环境的温度下降，我们把这样的反应或过程叫作\textbf{吸热}。

所以“冰袋变凉”这句话其实说反了：凉的是你的手和袋子表面，因为它们的热量被袋子里正在发生的过程夺走了。冰袋并没有“制造冷”，它是在\textbf{吞热}。

生活里这样的例子到处都是：

\begin{itemize}[itemsep=2pt]
\item 天然气灶上，甲烷燃烧剧烈放热，把锅里的水烧开；
\item 生石灰（\chem{CaO}）遇水放热，自热火锅靠它把菜汤加热到沸腾；
\item 酸碱中和放热，把稀盐酸倒进氢氧化钠溶液，烧杯会明显变温；
\item 小苏打遇上白醋会吸热，杯壁摸起来凉丝丝的（等下你就亲手试）；
\item 硝酸铵溶于水强烈吸热，这就是冰袋的秘密。
\end{itemize}

请注意最后一个：它甚至不是一个典型的“化学反应”，只是\textbf{溶解}。可见放热吸热与“是否生成了新物质”无关——只要粒子间的连接方式变了，能量账就可能不平。第 25 章讲过，溶解是粒子重新安排的过程；这一章我们要给这种“重新安排”标上价格。

溶解热还是双向的。硝酸铵溶于水吸热，可换一批溶质，账就反过来了：氢氧化钠溶于水明显放热——厨房管道疏通剂靠它腐蚀油脂，倒进水里你会感觉瓶子发烫；浓硫酸兑水更是剧烈放热，能把水溅开烫伤人（这是只在课堂上由老师演示的内容，家里绝对不许碰；而且规矩永远是“酸入水”，不能反过来）。同样是“把一种东西倒进水里”，有的变凉、有的滚烫、有的几乎没感觉——差别全在那两笔收支的差额上，等讲到冰袋时我们再把这笔账拆开。

还有一种吸热反应的“名场面”也值得一提：把氢氧化钡晶体和氯化铵晶体在烧杯里搅拌，烧杯会冷到把底下垫的湿木板\textbf{冻住}——杯子和木板之间那层水直接结冰，拿起烧杯能把木板一起带起来。两种固体搅一搅就造出“冰”，没有任何冰箱参与。这个实验涉及有毒的钡盐，只能看老师演示或视频，但它是“反应吸热”最有戏剧性的证据：环境的热被反应硬生生夺走，连水都被冻成了冰。

\section{把镜头推进到化学键}

价格标签贴在哪儿？贴在化学键上。

回忆第 17 到 19 章：化学键不是原子之间的小棍，而是让原子结合在一起时\textbf{能量更低的排布方式}。这句话反过来说就是：想把成键的原子拉开，你必须\textbf{输入}能量；而当原子结合成键，能量会\textbf{放出}。

可以拿磁铁打比方。两块吸在一起的磁铁，你要用力才能掰开——掰开的过程你在做功，能量存进了“分开”这个状态里；一松手，它们“啪”地吸回去，还会震你的手一下——吸合的过程能量放了出来。化学键也一样，只是把“磁力”换成了原子与电子之间的电性作用：

\begin{itemize}[itemsep=2pt]
\item \textbf{断键吸能}：拆开一根化学键，必须付出能量，从来不会白捡；
\item \textbf{成键放能}：形成一根化学键，一定放出能量，从来不会倒贴。
\end{itemize}

一场化学反应，按第 26 章的账本看是原子重新排队，按能量看则是两笔流水：先花钱拆掉旧键，再收钱建立新键。反应放热还是吸热，取决于\textbf{两笔钱的差额}：

\[ \text{净能量变化} \approx \text{断键吸收的能量} - \text{成键放出的能量} \]

差出来是负的（收的比花的多），系统能量下降，反应放热；差出来是正的，系统能量上升，反应吸热。就这么简单。放热反应的本质不是“释放了什么储存的热”，而是\textbf{产物里的键比反应物里的键更强、更“划算”}——原子们重新组合之后，整个体系落在了能量更低的位置。

\section{算一笔真实的账}

这根“价格标签”是可以查到数字的。化学家测量了各种化学键的\textbf{键能}——在气态下拆开 1 摩尔（\ud{6.02\times 10^{23}}{}根）某种键平均需要多少千焦。挑几根常见的：

\begin{table}[h]
\centering
\begin{tabular}{cc}
\toprule
化学键 & 键能（约，\ud{}{kJ\cdot mol^{-1}}） \\
\midrule
\chem{H-H} & 436 \\
\chem{O=O} & 498 \\
\chem{H-O} & 463 \\
\chem{C-H} & 413 \\
\bottomrule
\end{tabular}
\end{table}

现在来算氢气燃烧这笔账。第 26 章配平好的方程式是：

\[ \chem{2H_2} + \chem{O_2} \rightarrow \chem{2H_2O} \]

拆开旧键：2 根 \chem{H-H} 键加 1 根 \chem{O=O} 键，要付 $2\times 436+498=1370$ 千焦。建立新键：2 个水分子共 4 根 \chem{H-O} 键，能收回 $4\times 463=1852$ 千焦。收支相抵：$1370-1852=-482$ 千焦——每烧掉 2 摩尔氢气，体系净放出约 482 千焦的能量。

482 千焦是什么概念？够把大约 1.4 千克的水从室温烧到沸腾。而 2 摩尔氢气只有 4 克——小小一捧气体，藏着烧开一锅汤的能量。这就是火箭愿意背着液氢液氧上天的原因：按重量算，氢是放热效率最高的燃料之一。

\begin{qiaoliang}
大多数化学反应不在密封钢瓶里发生，而是在敞口的烧杯、锅碗和空气里——也就是在大气压基本不变的条件下。这时有个小麻烦：反应如果放热或生成气体，系统会膨胀，得花一点力气把周围的空气“推开”，这部分能量以“体积功”的形式溜走了，不等于你测到的热。化学家的处理办法很务实：干脆定义一个新的记账量\textbf{焓}（enthalpy，符号 $H$），把系统内部的能量加上“压强×体积”这一项打包。这样在\textbf{定压、不做其他功}的条件下，系统与环境交换的热量就恰好等于焓的变化：
\[ q_p = \Delta H \]
$\Delta H$ 读作“焓变”，等于产物的总焓减去反应物的总焓。约定从\textbf{系统的视角}记账：$\Delta H<0$，系统焓下降，能量流向环境，\textbf{放热}；$\Delta H>0$，系统焓上升，从环境吸走热量，\textbf{吸热}。你不需要会算 $H$ 的绝对值——化学家也只关心\textbf{变化量}，就像海拔可以随便设零点，高差才有意义。
\end{qiaoliang}

于是第 26 章的账本可以添上一栏。把焓变写在方程式后面，就叫\textbf{热化学方程式}：

\[ \chem{2H_2}(g) + \chem{O_2}(g) \rightarrow \chem{2H_2O}(g) \qquad \Delta H = -484\ \ud{}{kJ} \]

负号读作“放热”；括号里的 $g$ 是第 26 章介绍过的状态符号（气态），必须写上——因为同样的反应生成液态水和生成水蒸气，放出的热量是不一样的（蒸发本身要吸热，这是第 22 章分子间作用力的旧账）。

用文字把能量图画出来是这样的：纵轴是能量，反应物在高处，产物在低处，中间的落差就是 $\Delta H$。放热反应像水从高处流到低处；吸热反应则相反，产物坐在更高的台阶上，差额由环境埋单。

\begin{shiyan}
\textbf{亲手做一个“冰袋”：小苏打与白醋}

\textbf{材料}：小苏打（碳酸氢钠，厨房发面用的那种）两勺、白醋半杯、两个一次性杯子、一支温度计（没有就洗干净手用手指）、一根筷子。

\textbf{步骤}：① 白醋倒进杯子，先测（或摸）一下温度；② 一次倒入两勺小苏打，用筷子轻轻搅动；③ 边反应边观察温度变化，并用手背贴住杯壁。

\textbf{观察点}：剧烈冒泡的同时，温度计的读数是升还是降？杯壁是温还是凉？气体全部跑掉、反应停止之后，把杯子放着不动，温度又会怎样变化？

这个反应（碳酸氢钠遇酸生成二氧化碳、水和醋酸盐）总体是\textbf{吸热}的：它会从醋溶液、杯子和你的手里抢热量，所以杯壁发凉。反应停止后杯子慢慢回暖——那是环境在把热量还回来，系统与环境最终趋于同一温度。

\textbf{安全提示}：材料都可入口，很安全。但仍有两点：醋溅进眼睛会刺痛，搅动时悠着点；\textbf{绝对不要}把正在冒泡的混合物装进密闭瓶子摇晃——二氧化碳会把瓶子变成小炸弹。做完的液体可以直接倒进水槽冲掉。
\end{shiyan}

\section{能量账只看起点和终点}

现在来认识本章最省事、也最深刻的一条规律。

想象你从山脚下的小镇出发去山顶。可以走盘山公路，可以爬陡坡，可以先下到一个山谷再绕上去——不管你走哪条路、中途停下来吃了几顿饭，你的\textbf{海拔净升高}都一样，只取决于起点和终点。

焓就是这样的量。物理学家管它叫\textbf{状态函数}：它的变化只取决于系统的初态和末态，与中间路径无关。1840 年，俄国化学家\textbf{赫斯}（Germain Hess）在大量量热实验中总结出这条规律，今天叫\textbf{赫斯定律}：\textbf{一个反应无论一步完成还是分几步完成，总焓变相等}。

这条定律立刻显露出威力：有些反应的热根本没法直接测。比如碳不完全燃烧生成一氧化碳——你在实验里很难让碳“只烧一半”又烧得干干净净，副产物一堆。但可以绕路：

\begin{itemize}[itemsep=2pt]
\item 路线甲（一步）：碳完全燃烧，$\chem{C} + \chem{O_2} \rightarrow \chem{CO_2}$，实测 $\Delta H=-394$ 千焦；
\item 路线乙分两步：先 $\chem{C} + \frac{1}{2}\chem{O_2} \rightarrow \chem{CO}$（设为 $x$，测不准），再 $\chem{CO} + \frac{1}{2}\chem{O_2} \rightarrow \chem{CO_2}$，实测 $-283$ 千焦。
\end{itemize}

两条路线起点相同（碳加足量氧）、终点相同（二氧化碳），总账必须相等：$x+(-283)=-394$，解出 $x=-111$ 千焦。一步测不出的反应热，用两把能测的“尺子”拼出来了。碳变金刚石、葡萄糖发酵、各种又慢又危险的反应——它们的焓变很多都是这样“绕道”算出来的，根本不需要真把反应做一遍。

这套“纸上绕路”的功夫，是现代化学工业的隐形地基。设计一座合成氨工厂或硫酸工厂时，工程师面对的核心难题之一就是\textbf{热量管理}：反应放出的热散不掉，设备会过热失控；该补的热跟不上，反应又会“熄火”。他们的做法是先把整条工艺拆成一步一步的反应，从手册和量热数据里查出每一步的焓变，再用赫斯定律把它们加加减减，算出整条生产线每一段的净吸放热，据此设计换热器和冷却系统——在图纸阶段就把能量账算平，而不是等工厂爆炸了再总结教训。至于为什么工厂还要费尽心机控制温度、请催化剂帮忙，那是第 30、31 章的故事。

\begin{zhengju}
你可能会觉得赫斯定律“显然如此”——能量守恒嘛，起点终点定了，差额还能跑掉？但在 1840 年，这一点都不显然。当时占上风的还是“热质说”：热被想象成一种没有重量、无孔不入的流体，叫“热质”，物体含热质多就热，少就冷。按这个图景，热是一种\textbf{东西}，反应放热就是“释放热质”，没有道理非跟路径无关不可。

赫斯的底气来自实验。他用自己改进的量热计测量酸碱中和、盐类溶解的热效应，发现不管把反应拆开走几步，总热量都对得上账。他不满足于记录现象，而是把这个“总账守恒”提升为普遍规律——这在热质说框架里是无法解释的怪事，却在几十年后能量守恒定律确立时，成为它最漂亮的注脚之一：热不是物质，而是\textbf{能量的一种转移形式}，能量只会转化，不会凭空增减。

给“热质说”棺材钉上钉子的实验也很精彩。1798 年，伦福德伯爵在兵工厂监督给大炮钻孔，发现钻头和大炮摩擦生热似乎“用不完”——只要一直钻，热就一直冒出来。如果热是藏在金属里的“热质”，钻下来的碎屑总该把热质放完吧？可碎屑带走的“热质”少得可怜，钝钻头几乎不钻下任何东西，却照样大量生热。结论：热可以靠做功\textbf{源源不断地产生}，它不是库存的货物，而是能量的流动。

而焓变的测量工具也在同期进化。早在 1780 年代，拉瓦锡和拉普拉斯就造出了“冰量热计”：把反应（甚至一只豚鼠！）放在冰层包裹的腔室里，称一称融化了多少冰，就知道放出了多少热。拉瓦锡用这台机器证明：豚鼠“取暖”放出的热，和它呼吸消耗的氧气，与木炭燃烧遵循同一本账——\textbf{生命的热和炉火的热是同一种东西}。这是化学通向生物学的第一座桥，我们这本书后半程会一直走在它上面。
\end{zhengju}
\section{反应热是怎么测出来的：量热法}

赫斯的数据不是从天上掉下来的，是从一台台\textbf{量热计}里读出来的。量热的原理朴素得可爱：让反应在一个热量跑不掉的“小房间”里发生，看房间的温度变了多少。

最简单的版本你在学校实验室就能搭：一个保温杯（泡沫塑料杯也行）、一支温度计、一根搅拌棒。在杯里做反应——比如酸碱中和——测温度变化，然后算账。水有一个特别好记的脾气：让 1 克水升高 $1\,^{\circ}\mathrm{C}$，需要约 4.2 焦的热量（这叫水的比热容）。于是：

\[ q = m \times 4.2\ \ud{}{J/(g\cdot ^{\circ}C)} \times \Delta T \]

100 克溶液升温 $5.4\,^{\circ}\mathrm{C}$，就吸收了约 $100\times4.2\times5.4\approx2270$ 焦——这部分热只能是反应放出来的。除以反应物的物质的量，就得到每摩尔的焓变。稀盐酸中和氢氧化钠，这样测出来约是 $-57$ 千焦每摩尔：每生成 1 摩尔水，放出 57 千焦。

测食物的热量就要请更结实的设备了：\textbf{弹式量热计}。把干燥的食物样品塞进密封耐压的“钢弹”，充入纯氧，点火，让食物彻底烧光；钢弹泡在定量的水里，水温升多少，食物就含多少化学能。你吃的饼干包装上印的“能量 2000 千焦”，就是这么烧出来的数字。有意思的是，同一块饼干在你身体里不是“烧掉”的，而是被酶领着走几十步温和的反应（第三册会细讲）——但按赫斯定律，起点（饼干加氧气）和终点（二氧化碳加水）相同，总账分毫不差。你的身体，就是一座走盘山公路的弹式量热计。

\section{这套模型的边界在哪里}

焓变很强大，但有几件事它管不了，本章先把丑话说在前面：

\begin{itemize}[itemsep=2pt]
\item \textbf{$\Delta H$ 不等于反应的方向。}放热只是“能量账盈余”，可冰袋里的硝酸铵明明在吸热，为什么不用插电它就自己溶解了？可见除了能量高低，自然界还另有一本“分散程度”的账。这是第 28 章的主角——熵。
\item \textbf{$\Delta H$ 不等于反应的快慢。}氢气和氧气混合在一起，账面上早该放出 484 千焦，可它们能相安无事地放好几年，直到一颗火星引爆全场。能量落差决定“值不值得”，启动门槛决定“动不动得了”。门槛的事，第 29 章讲。
\item \textbf{焓变是有条件的数字。}它默认定压；温度、压强、物质状态变了，数值也会跟着变。手册上查到的通常是 $25\,^{\circ}\mathrm{C}$、1 个大气压、各物质处于最稳定状态的“标准焓变”，换个场合要修正。
\item \textbf{键能是平均值的估算。}前面那张键能表是统计平均——“一根 \chem{H-O} 键 463 千焦”，但水分子里第一根 \chem{H-O} 和第二根断起来其实不一样费劲。用键能估焓变，误差几个百分点很正常；精确数字要靠量热计实测。
\end{itemize}

\begin{wuqu}
\textbf{“化学键里储存着能量，断开化学键会释放能量。”}——这可能是化学里流传最广的一句错话，连一些科普读物都在重复。

反例就在本章的键能表里：拆开 1 摩尔 \chem{H-H} 键，必须\textbf{输入} 436 千焦，一分不少。如果断键放能，氢气自己就该散成原子、还顺便发电，宇宙早就不是现在这个样子了。

打两个比方帮你拨乱反正。其一：化学键像两块吸住的磁铁，掰开要费劲（吸能），吸合才“啪”地放能——能量藏在“分开”这个状态里，而不是藏在“连接”里。其二：说“键里储存能量”，就像说“山谷底部储存着高度”——恰恰相反，能量储存在高处。氢气燃烧放热，不是因为“拆键释放能量”，而是因为 \chem{H-O} 键比 \chem{H-H} 和 \chem{O=O} 键更强，成新键收回来的钱多于拆旧键付出去的钱。\textbf{放热的永远是净差额，不是断键本身。}（这个误解后面还会换个马甲出现——比如“ATP 的高能键一断就释放能量”，第三册再拆它一次。）
\end{wuqu}

\section{回到那两只小袋子}

现在可以兑现开场的承诺了。

\textbf{暖宝宝}里装的是铁粉，外加活性炭、食盐、水和保温的蛭石。撕开塑料袋，空气涌进去，铁开始生锈：

\[ \chem{4Fe} + \chem{3O_2} \rightarrow \chem{2Fe_2O_3} \qquad \Delta H \approx -1650\ \ud{}{kJ} \]

这笔账大得惊人：每氧化 4 摩尔（约 224 克）铁，放出约 1650 千焦。一片暖宝宝约含 30～40 克铁粉，全部氧化能放出二百多千焦——够把两升水烧到快开。之所以只觉“温温的”，是因为食盐和水把反应调成慢速档，蛭石把热量兜住慢慢放，几小时才放完这笔巨款。所以暖宝宝的热不是“藏”在袋子里的，而是\textbf{铁与氧之间更强的键}换来的——撕袋子之前，反应物就在，只是被塑料袋挡着见不了面。

\textbf{冰袋}里是硝酸铵颗粒和一小袋水。捏破内袋，\chem{NH_4NO_3} 开始溶解。溶解也有能量账：把离子从整齐的晶体里拆散，要花钱（破坏第 18 章的离子键，吸能）；离子被水分子团团围住形成水合离子，能收钱（第 25 章的溶剂化，放能）。对大多数盐（比如食盐）两笔账大致相抵，水温几乎不变；但硝酸铵是特例——拆晶格花的钱明显多过水合收的钱，净吸热约 $+26$ 千焦每摩尔。差额从哪儿来？从溶液、从袋子、从你扭伤的脚踝里抽。你的脚踝变凉，是因为热量被借去填了溶解的能量缺口。

一热一冷，同一套账本：放热还是吸热，永远是“拆旧连接”与“建新连接”的差额，没有例外。

顺便说一句这两只袋子的“职业身份”：它们都是正经的医疗和运动用品。冷敷让扭伤部位的血管收缩、减缓肿胀，热敷促进血液流动、缓解肌肉僵硬——原理在医学上都有依据，但请留意本书反复强调的规矩：敷多久、多高的温度、什么人适合，要看说明书、听医嘱，冻伤和低温烫伤都是真实存在的风险。化学反应只管放热吸热，怎么用是人要负责的部分。

\section{能量账之外，还有一本账}

不过，细心的你应该已经发现一个漏洞：既然吸热意味着系统能量上升，硝酸铵为什么还“抢着”溶解？把手伸进冰袋，溶解会自己停下来吗？不会。它一路溶解到饱和为止，好像有什么东西在推着它往能量\textbf{更高}的地方走。

这正是本章模型管不了、却无比重要的那部分：能量下降不是自然界唯一的偏好。粒子和能量还偏爱“分散开”——晶体里整齐排列的离子，一旦有机会散进水里，常常就散开了，哪怕为此付出吸热的代价。这本“分散的账”叫熵，它和焓加在一起，才真正决定一个过程\textbf{朝哪个方向走}。下一章，我们就去算这第二本账。

\begin{tiaozhan}
\textbf{焓和内能差多少？}（跳过不影响主线）化学家还定义一个更“内部”的量：内能 $U$，只算系统粒子的动能和相互作用能，不管推空气的账。定压下 $\Delta H=\Delta U+\Delta(pV)$；对气体反应，用理想气体近似可写成 $\Delta H=\Delta U+\Delta n_{\text{气}}RT$，其中 $\Delta n_{\text{气}}$ 是气体摩尔数的变化，$R$ 是气体常数，$T$ 是温度。拿氢气燃烧（生成气态水）来说：$\Delta n_{\text{气}}=2-3=-1$，在 $25\,^{\circ}\mathrm{C}$ 时 $\Delta n_{\text{气}}RT\approx-2.5$ 千焦——焓变与内能变只差约 0.5\%。对不放出也不消耗气体的反应（比如溶液里的中和），两者几乎相等。这就是为什么化学家偏爱焓：它把“推开大气”这件琐事默默处理掉了。
\end{tiaozhan}

\section*{练一练}

\begin{sikaoti}
\item 画一幅能量示意图：纵轴标“能量”，画出甲烷燃烧（放热）和硝酸铵溶解（吸热）两条线，各标出反应物、产物的位置和 $\Delta H$ 的正负。图旁注明：这是表示能量关系的示意图，不是粒子的真实照片。
\item 用本章键能表估算：1 摩尔甲烷（\chem{CH_4}，含 4 根 \chem{C-H} 键）全部拆成气态原子需要约多少能量？再查资料找到实测值，想一想为什么和你估的不完全一样。
\item 某同学用保温杯做中和实验：100 克稀盐酸与 100 克氢氧化钠溶液混合，温度从 $25.0\,^{\circ}\mathrm{C}$ 升到 $28.2\,^{\circ}\mathrm{C}$。估算这次反应放出了多少热量（溶液的比热容近似取水的 4.2 \ud{}{J/(g\cdot ^{\circ}C)}）。如果改用 200 克溶液做同样反应，温度升高会怎么变？
\item 暖宝宝用完后变硬、变重。从粒子角度解释：增加的重量是从哪里来的？为什么这与质量守恒不矛盾？
\item “放热反应不需要加热就能发生。”这句话对吗？请用本章内容和氢氧混合气体的例子各给一个反驳，并说明“放热”“自发”“快速”是三件不同的事。
\item 设计一个实验，比较等量食盐和硝酸铵溶于水时的温度变化。写出你要控制哪些变量、预测会看到什么，并解释两种盐的差别本质上差在哪两笔“收支”上。
\end{sikaoti}
